% =====================================================================
%  CAPÍTULO 7 · CONCLUSIONES Y TRABAJO FUTURO
%  Función: cerrar el círculo. Cada objetivo de la Sección 1.3 se retoma
%  UNO A UNO: cumplido / parcialmente / no cumplido, con su evidencia.
%  Reglas clave:
%   - No introduzcas aquí resultados nuevos que no estén en el Cap. 6.
%   - Limitaciones (lo que no salió) separadas del trabajo futuro.
%   - Trabajo futuro: ~1 página, concreto y ligado a limitaciones reales.
% =====================================================================
\chapter{Conclusiones y trabajo futuro}
\label{cap:conclusiones}

Este capítulo valora el grado de cumplimiento de los objetivos, resume las
aportaciones y las limitaciones del trabajo y propone líneas de continuación.

\section{Cumplimiento de los objetivos}
\label{sec:cumplimiento}

La Tabla~\ref{tab:cumplimiento} resume el estado final de cada objetivo
específico planteado en la Sección~\ref{sec:objetivos}.

\begin{table}[htbp]
  \centering
  \caption{Grado de cumplimiento de los objetivos.}
  \label{tab:cumplimiento}
  \small
  \begin{tabularx}{\textwidth}{@{}lLlL@{}}
    \toprule
    \textbf{Obj.} & \textbf{Descripción} & \textbf{Estado} & \textbf{Evidencia} \\
    \midrule
    \ref{obj:estado}      & Comparativa de soluciones   & Cumplido & Sección~\ref{sec:soluciones} \\
    \ref{obj:app}         & Aplicación Android funcional & Cumplido & Tabla~\ref{tab:casos} \\
    \ref{obj:rendimiento} & Latencia $<$ \qty{100}{ms}   & Cumplido & Figura~\ref{fig:latencia} \\
    \ref{obj:validacion}  & Validación con usuarios      & Parcial  & SUS = 81, pero $n = 8$ (Sección~\ref{sec:discusion}) \\
    \bottomrule
  \end{tabularx}
  \fuente{Elaboración propia.}
\end{table}

Desarrolla después cada objetivo en uno o dos párrafos. Por ejemplo: el
objetivo~\ref{obj:rendimiento} se ha cumplido, ya que con la caché la latencia
media se mantiene en \qty{92}{ms} con 400 usuarios, el doble de la carga exigida
por RNF-01. El objetivo~\ref{obj:validacion} se considera parcialmente cumplido:
la puntuación supera el umbral, pero la muestra es pequeña.

\section{Conclusiones}
\label{sec:conclusiones}

Resume qué aporta el trabajo y qué se ha aprendido, tanto en lo técnico como en
la gestión del proyecto. Las conclusiones se apoyan en los resultados del
Capítulo~\ref{cap:pruebas}; no repitas la introducción ni añadas afirmaciones
que no se hayan demostrado.

\section{Limitaciones}
\label{sec:limitaciones}

Enumera con honestidad lo que el trabajo no cubre o no ha podido validar y por
qué: por ejemplo, la ausencia de pruebas en varios servidores o la muestra
reducida de la prueba de usabilidad.

\section{Trabajo futuro}
\label{sec:futuro}

Propón tres o cuatro líneas concretas y priorizadas, cada una ligada a una
limitación. Por ejemplo: repetir la prueba de usabilidad con al menos treinta
usuarios de distintos grados; desplegar el servidor en dos nodos para medir la
escalabilidad horizontal, y publicar una versión para iOS.
